\documentclass[11pt]{article}
\usepackage[margin=25mm]{geometry}
\usepackage{amsmath,amssymb,hyperref}
\title{Exact finite-target seeds and quotient evaluation for a Coxeter search}
\author{Bounded L033 method component}\date{}
\begin{document}\maketitle
\begin{abstract}
We replace two operations in the retained LINS normal-subgroup search:
finite-target epimorphism enumeration and evaluation of a composed
homomorphism. Both changes preserve the complete set of normal kernels.
No new hyperbolic quotient, minimum, named group, or completed target bound
is asserted by this component.
\end{abstract}

\section{The finite-target list}
Let
\[
 W=\langle s_1,\ldots,s_r\mid s_i^2=1,
       (s_i s_j)^{m_{ij}}=1\ (i<j)\rangle
\]
for a specified finite Coxeter matrix. Our interface checks the literal
presentation, including every relator, against this matrix. Fix a finite
target group $Q$. Every homomorphism $W\to Q$ corresponds to a tuple
$(a_1,\ldots,a_r)$ with $a_i^2=1$ and
$(a_i a_j)^{m_{ij}}=1$. Identities and proper divisors of the prescribed
orders are allowed at this stage. Surjectivity is exactly
$\langle a_1,\ldots,a_r\rangle=Q$.

After choosing a prefix $(a_1,\ldots,a_{k-1})$, its stabilizer under
simultaneous inner conjugation is
$C_Q(a_1,\ldots,a_{k-1})$. The possible next elements are exactly the
square-one elements satisfying the relations to all chosen marks. This
set is invariant under that stabilizer. Choose one representative from
each of its actual conjugation orbits and continue with the centralizer
of the extended prefix. Induction proves that the resulting full tuples
represent every inner-conjugacy orbit exactly once. Filtering the full
tuples by generation and constructing the actual homomorphism therefore
finds every epimorphism up to inner conjugation. Inner-conjugate
homomorphisms have the same kernel. Outer-conjugate ones may still have
the same kernel; the normal-subgroup graph removes these duplicates.
No quotient-isomorphism shortcut is used to remove distinct kernels.

The retained default LINS routine first seeds its graph with all kernels
of epimorphisms to the finite targets in its table through the bound.
Its installed \texttt{findTQuotients.gi} explicitly uses
\texttt{GQuotients} for this operation. Our initialization inserts
precisely the same set of kernels through the existing graph-insertion
interface. It then disables only that already completed target-seeding
pass. All subsequent prime-power quotient, intersection, and graph
operations are unchanged, except the evaluation below. The invocation
uses the default options apart from this initialization and target-list
filter; no cutting or early termination option is supplied. The caller
retains the original index guard. The LINS algorithm and its finite
classification tables, through their installed $10^7$ limit, remain
explicit computational and classification premises, as in the accepted
eight-diagram core. We do not reprove their completeness.

\section{Evaluation in the actual quotient}
At a prime-power step LINS constructs an isomorphism
$\iota:H\to\widehat H$, where $\widehat H$ is finitely presented, and an
epimorphism $\mu:\widehat H\to E$ onto an elementary abelian group.
It obtains the conjugation matrix by evaluating words
$w=xyx^{-1}\in H$ as $\mu(\iota(w))$.
Set $\nu=\mu\circ\iota$. Then $\nu(w)=\mu(\iota(w))$ for every such
word. Consequently every matrix entry, every invariant submodule, every
pullback kernel, and every recursive search remains identical.

The process-local override adds the local variable \texttt{Nu}, sets
\texttt{Nu:=Iso*Mu} in GAP's composition notation, and replaces
only that nested image by \texttt{Image(Nu,w)}. All other executable
statements are literal copies of the pinned installed function. No
installed package is edited. GAP's \texttt{ghomfp.gi} composition method
copies the augmented coset table and maps its generator images through
the second map. Rewriting can thus evaluate directly in $E$ instead of
building a long intermediate word in $\widehat H$. This implementation
explains the possible saving; the mathematical equality above does not
depend on a timing assumption or an injectivity claim for $\mu$.

\paragraph{Controls and scope.}
A literal traversal of all square-one tuples independently agrees with
the inner-orbit traversal for three matrices and six finite groups
(18 cases, seven surjective inner orbits). Complete old and new normal
kernel inventories agree bijectively as ordered marked quotients for
three bounded presentations. During those changed searches all 676
composed conjugation images are also evaluated by the original nested
formula and compared exactly. Invalid matrices and mismatched source
presentations are rejected. The independent core checker still requires
full intersection and actual flag-transitivity tests for positive
hypertopes; this method does not weaken them.

The preserved first control driver accidentally captured the replacement
function as its reference. Its result is not used for the original-versus-new
traversal comparison; the corrected driver reruns all controls while
capturing the installed function first. A pair of capped diagnostic
profiles shows one index-1920, dimension-71 matrix step taking about
12.7 seconds by nested evaluation and 1.1 seconds by composition.
Neither profile completes its target bound, and neither is a new census
or general performance theorem. All earlier capped searches remain
uncredited.

\paragraph{Attribution.}
The finite targets and surrounding algorithm are the installed LINS 0.9
implementation, based on Firth's normal-subgroup algorithm; the direct
composition method is already present in GAP 4.15.1. Stabilizer-orbit
traversal is standard. The present author supplies the Coxeter interface,
exact substitution proof, code, and bounded controls. The root supplied
the original L033 probe; a separate reader supplied the acquired LINS
interface audit. No novelty, whole-source, measurement, question, or Lean
credit is claimed.
\end{document}
